\chapter{गर्भनिरोध}\label{ch:g8:contraception}

पिछले अध्यायों ने ग़ज़ब की ताक़त वाली एक कार्य-प्रणाली का नक़्शा बनाया:
दो \omterm{def:g5:first-look-at-cells:cell}{कोशिकाएँ} मिलती हैं, और चालीस हफ़्ते बाद एक इंसान होता है। मनुष्य ही
अकेली ऐसी \omterm{def:g5:biodiversity:species}{प्रजाति} है जो अपनी कार्य-प्रणाली समझती है --- और समझ के साथ
वह चुनाव भी आता है जिसका सामना मेंढक को कभी नहीं करना पड़ता: उसे चलाना
है या नहीं, और कब। \emph{\omterm{prop:g8:contraception:principle}{गर्भनिरोध}} उसी चुनाव का व्यावहारिक विज्ञान है,
और उसका हर तरीक़ा किसी ऐसे क़दम को रोककर काम करता है जिसका नाम अब तुम
बता सकते हो।

\section{उसूल: क़दम का नाम लो, और उसी क़दम को रोको}

\begin{proposition}[गर्भनिरोध कार्य-प्रणाली पर ही बैठता है]\label{prop:g8:contraception:principle}
\omterm{def:g5:human-reproduction:pregnancy}{गर्भावस्था} के लिए एक अटूट कड़ी चाहिए
(\cref{prop:g8:sexual-reproduction:core},
\cref{prop:g8:from-egg-to-birth:firstweek}): \omterm{def:g8:sexual-reproduction:gametes}{युग्मक} बनना --- शुक्राणुओं
का पहुँचाया जाना --- \omterm{def:g8:reproductive-systems:female}{अंडवाहिनी} तक की तैराकी --- किसी \omterm{def:g4:animal-reproduction:sexual}{अंडाणु} का छोड़ा
जाना और मिलना --- \omterm{def:g5:flower-to-fruit:steps}{निषेचन} --- आरोपण। \emph{गर्भनिरोध}\index{गर्भनिरोध}
इस कड़ी का कोई भी जान-बूझकर किया गया टूटना है। दवाख़ाने की अलमारी का हर
तरीक़ा किसी एक नामज़द कड़ी पर हमला है, और उन्हें उनकी कड़ी के हिसाब से
छाँट लेना ही पूरा सिद्धांत है:
\begin{itemize}
  \item \emph{रोक} वाले तरीक़े \omterm{def:g8:sexual-reproduction:gametes}{युग्मकों} को मिलने से रोकते हैं;
  \item \emph{\omterm{def:g8:puberty:hormones}{हॉर्मोन}} वाले तरीक़े \omterm{def:g4:animal-reproduction:sexual}{अंडाणु} का छोड़ा जाना ही रोक देते
        हैं;
  \item और बाक़ी तरीक़े मिलन या आरोपण को नाकाम कर देते हैं।
\end{itemize}
\end{proposition}
\admitted

\section{मुख्य तरीक़े, उनकी कड़ी के हिसाब से}

\begin{definition}[रोक वाले तरीक़े: कंडोम]\label{def:g8:contraception:condom}
\emph{कंडोम}\index{कंडोम} --- यानी \omterm{def:g8:reproductive-systems:male}{शिश्न} पर पहना जाने वाला एक पतला ग़िलाफ़
(एक \omterm{def:g4:animal-reproduction:sexual}{मादा} रूप योनि पर अस्तर बनता है) --- कड़ी को पहुँचाने वाले क़दम पर
रोक देता है: वीर्य साथी के शरीर में जाता ही नहीं। यह हर जगह मिलता है,
और सिर्फ़ ज़रूरत के वक़्त चाहिए --- और अलमारी में अकेला है एक दूसरी सेवा
के लिए भी: यह इकलौता तरीक़ा है जो सहवास से फैलने वाले \emph{संक्रमणों}
का रास्ता भी रोकता है, यानी \cref{ch:g9:microbes-and-infection} वाली
दुनिया का विषय। यही दोहरी ड्यूटी है जिसकी वजह से कोई और तरीक़ा चल रहा हो
तब भी कंडोम तस्वीर में बना रहता है।
\end{definition}

\begin{definition}[हॉर्मोन वाले तरीक़े: गोली]\label{def:g8:contraception:pill}
\emph{गर्भनिरोधक गोली}\index{गर्भनिरोधक गोली}
\cref{def:g8:puberty:hormones} की अपनी ही मशीनरी इस्तेमाल करती है: रोज़
ली जाए तो उसके \omterm{def:g8:puberty:hormones}{हॉर्मोन} \omterm{def:g8:reproductive-systems:female}{अंडाशयों} के चक्र को उसी ठहरी हुई हालत में थामे
रखते हैं (\cref{prop:g8:reproductive-systems:cycle} की चौथी घटना, पर
बिना किसी \omterm{def:g5:human-reproduction:pregnancy}{गर्भावस्था} के) --- \emph{न \omterm{prop:g8:reproductive-systems:cycle}{अंडोत्सर्ग}, न \omterm{def:g4:animal-reproduction:sexual}{अंडाणु}, और निषेचित
करने को कुछ भी नहीं}। यह डॉक्टर के नुस्ख़े से मिलती है, जो तरीक़े को
व्यक्ति से मिलाते हैं; नियम से ली जाए तो बेहद भरोसेमंद --- और संक्रमणों
के ख़िलाफ़ कुछ भी नहीं देती, इसीलिए \omterm{def:g8:contraception:condom}{कंडोम} की वह पक्की साझेदारी।
\end{definition}

\begin{example}[अलमारी, कड़ी के हिसाब से छँटी हुई]\label{ex:g8:contraception:shelf}
बाक़ी आम तरीक़े, उनकी हमला की गई कड़ी के हिसाब से दर्ज: \omterm{def:g8:puberty:hormones}{हॉर्मोन} वाले
प्रत्यारोपण और चिपकाऊ पट्टियाँ --- यानी गोली का ही उसूल, पर बेहतर
याददाश्त के साथ (महीनों या सालों तक, बिना रोज़ लिए); वे छोटे उपकरण जो
डॉक्टर \omterm{prop:g5:human-reproduction:plan}{गर्भाशय} में रख देते हैं --- यानी मिलन और आरोपण, सालों तक एक साथ
नाकाम किए हुए; और \emph{आपात गोली}, यानी बाद में कुछ दिनों के भीतर ली
जाने वाली \omterm{def:g8:puberty:hormones}{हॉर्मोन} की एक ख़ुराक, जो मुख्य रूप से अभी न हुए \omterm{prop:g8:reproductive-systems:cycle}{अंडोत्सर्ग} को
टालकर काम करती है --- यानी आख़िरी उपाय वाला एक ब्रेक, जो दवाख़ाने पर
मिलता है, और जो \emph{पहले से} इस्तेमाल किए गए किसी तरीक़े की जगह हरगिज़
नहीं ले सकता। और इनमें से कोई भी जो नहीं करता: संक्रमण के ख़िलाफ़ कुछ
भी।
\end{example}

\begin{remark}[क्या काम नहीं करता]\label{rem:g8:contraception:notwork}
मैदान की गप वाले ये तरीक़े, इसी कड़ी के सामने नापे हुए: ``सुरक्षित दिन
गिनना'' तो \cref{rem:g8:reproductive-systems:living} वाली बदलाहट पर धराशायी हो जाता
है --- क्योंकि \omterm{prop:g8:reproductive-systems:cycle}{अंडोत्सर्ग} की तारीख़ पहले से जानी ही नहीं जा
सकती, और \omterm{def:g8:reproductive-systems:male}{शुक्राणु} तो उसका इंतज़ार करते हुए कई दिन तक जीते हैं; और ``वक़्त पर पीछे हट जाना'' भी इस
तथ्य पर धराशायी होता है कि \omterm{def:g8:reproductive-systems:male}{शुक्राणु} आख़िरी घड़ी से पहले ही सफ़र कर चुके
होते हैं। दोनों ठीक वहीं नाकाम होते हैं जहाँ नामज़द-कड़ी वाला सिद्धांत
कहता है: वे भरोसे से कोई कड़ी रोकते ही नहीं। गप कोई तरीक़ा नहीं है।
\end{remark}

\section{चुनाव, ज़िम्मेदारी, और मदद का नक़्शा}

\begin{proposition}[यह चुनाव असली है और साझा भी]\label{prop:g8:contraception:choice}
वे तथ्य, जो यह कार्य-प्रणाली थोप देती है:
\begin{itemize}
  \item \omterm{def:g5:human-reproduction:pregnancy}{गर्भावस्था} के लिए एक ही मौक़ा काफ़ी है, किसी रिश्ते की लंबाई
        नहीं: यह कड़ी मौक़े गिनती ही नहीं;
  \item यह एक \emph{साझी} कड़ी है: दोनों साथियों की \omterm{def:g5:first-look-at-cells:cell}{कोशिकाएँ}, दोनों
        की ज़िम्मेदारी --- \omterm{def:g8:contraception:condom}{कंडोम} और गोली किसी एक-एक लिंग का मामला
        नहीं, बल्कि एक बातचीत हैं;
  \item और परहेज़ के सिवा कोई भी तरीक़ा मुकम्मल नहीं है, हालाँकि अच्छे
        तरीक़े ठीक-ठीक वैसे इस्तेमाल हों जैसे वे बने हैं तो
        क़रीब-क़रीब मुकम्मल हैं --- और असल दुनिया की ज़्यादातर
        नाकामियाँ इसी ``ठीक से इस्तेमाल'' में बसती हैं।
\end{itemize}
\end{proposition}
\admitted

\begin{example}[मदद का नक़्शा]\label{ex:g8:contraception:help}
\Cref{ex:g8:puberty:questions} वाले पते, और आगे बढ़ाकर: स्कूल की नर्सें
और घर के डॉक्टर \omterm{prop:g8:contraception:principle}{गर्भनिरोध} के सवालों के जवाब गोपनीयता के साथ देते हैं ---
और नाबालिगों के भी; परिवार नियोजन केंद्र हर इलाक़े में ठीक इसी के लिए
हैं, मुफ़्त और बिना किसी फ़ैसले के; और दवाख़ाने वाले आपात गोली की घड़ी
सँभालते हैं। यह पूरा इंतज़ाम इसीलिए बना है कि किसी को यह कड़ी गप के सहारे
पार न करनी पड़े --- मदद का नक़्शा ख़ुद तरीक़े का हिस्सा है।
\end{example}

\begin{method}[अपने सुने हुए किसी भी तरीक़े को परखना]\label{met:g8:contraception:evaluate}
किसी भी दावे वाले तरीक़े के बारे में, क्रम से पूछो:
\begin{enumerate}
  \item वह कड़ी की \emph{कौन-सी कड़ी} रोकता है --- क्या तुम उसे नक़्शे
        पर नाम लेकर बता सकते हो?
  \item और वह \emph{कितने भरोसे} से रोकता है, जैसे असली लोग उसे
        इस्तेमाल करते हैं?
  \item क्या वह \emph{संक्रमणों} से भी बचाता है, या उसे \omterm{def:g8:contraception:condom}{कंडोम} की
        साझेदारी चाहिए?
  \item और वह माँगता क्या है --- रोज़ की याददाश्त, कोई नुस्ख़ा, कोई
        फ़िटिंग --- और क्या वह इस्तेमाल करने वाले पर बैठता है?
\end{enumerate}
जिसकी कोई कड़ी नाम लेकर न बताई जा सके, वह तरीक़ा नहीं
(\cref{rem:g8:contraception:notwork} पहले ही सवाल पर नाकाम हो जाती
है)।
\end{method}

\section{अभ्यास}

\begin{exercise}[$\star$]\label{exo:g8:contraception:1}
\omterm{def:g5:human-reproduction:pregnancy}{गर्भावस्था} के लिए ज़रूरी कड़ी बताओ, और उसी के सामने \omterm{prop:g8:contraception:principle}{गर्भनिरोध} की
परिभाषा दो।
\end{exercise}

\begin{exercise}[$\star$]\label{exo:g8:contraception:2}
\omterm{def:g8:contraception:condom}{कंडोम} कौन-सी कड़ी रोकता है? और गोली?
\end{exercise}

\begin{exercise}[$\star$]\label{exo:g8:contraception:3}
\omterm{def:g8:contraception:condom}{कंडोम} अकेला कौन-सी दूसरी सेवा देता है?
\end{exercise}

\begin{exercise}[$\star$]\label{exo:g8:contraception:4}
आपात गोली मुख्य रूप से कैसे काम करती है, और वह है क्या नहीं?
\end{exercise}

\begin{exercise}[$\star$]\label{exo:g8:contraception:5}
मदद के नक़्शे के तीन पते बताओ, और यह भी कि हर एक क्या देता है।
\end{exercise}

\begin{exercise}[$\star$]\label{exo:g8:contraception:6}
``सुरक्षित दिन गिनना'' भरोसे लायक़ क्यों नहीं है?
\cref{rem:g8:contraception:notwork} में से दो कारण दो।
\end{exercise}

\begin{exercise}[$\star\star$]\label{exo:g8:contraception:7}
\cref{ex:g8:contraception:shelf} वाली अलमारी को हमला की गई कड़ी के
हिसाब से छाँटो, और उसमें \omterm{def:g8:contraception:condom}{कंडोम} तथा गोली भी जोड़ो।
\end{exercise}

\begin{exercise}[$\star\star$]\label{exo:g8:contraception:8}
डॉक्टर नौजवान जोड़ों से इतनी बार ``गोली \emph{और} \omterm{def:g8:contraception:condom}{कंडोम}'' क्यों कहते
हैं? दो हिफ़ाज़तें, हर एक पर एक वाक्य।
\end{exercise}

\begin{exercise}[$\star\star$]\label{exo:g8:contraception:9}
\cref{prop:g8:reproductive-systems:cycle} की मदद से समझाओ कि गोली चक्र
को किस हालत में थामे रखती है, और वह हालत \omterm{def:g5:human-reproduction:pregnancy}{गर्भावस्था} को क्यों रोक देती
है।
\end{exercise}

\begin{exercise}[$\star\star$]\label{exo:g8:contraception:10}
``दोनों साथियों की \omterm{def:g5:first-look-at-cells:cell}{कोशिकाएँ}, दोनों साथियों का मामला।'' \omterm{prop:g8:contraception:principle}{गर्भनिरोध} को
जोड़े के किसी एक ही पक्ष पर छोड़ देने की आदत के ख़िलाफ़ इस वाक्य का बचाव
करो।
\end{exercise}

\begin{exercise}[$\star\star$]\label{exo:g8:contraception:11}
\cref{met:g8:contraception:evaluate} इन पर चलाओ: (क) \omterm{def:g8:puberty:hormones}{हॉर्मोन} वाला
प्रत्यारोपण; (ख) ``वक़्त पर पीछे हट जाना''।
\end{exercise}

\begin{exercise}[$\star\star\star$]\label{exo:g8:contraception:12}
``\omterm{prop:g8:contraception:principle}{गर्भनिरोध} व्यावहारिक शरीर-क्रिया विज्ञान है: अलमारी का हर तरीक़ा इसी
किताब का एक अध्याय है, निशाना साधे हुए।'' तीन तरीक़ों और उनके तीन
अध्यायों से इसे सही ठहराओ।
\end{exercise}

\section{समस्या: परिवार नियोजन का परामर्श}

\begin{problem}[{सप्ताहांत समस्या --- नियोजन केंद्र पर सवालों से भरी एक
दोपहर}]\label{pb:g8:contraception:1}
किसी परिवार नियोजन सलाहकार की एक (कल्पित) दोपहर: चार आगंतुक, चार
हालात, और दीवार पर कड़ी का एक नक़्शा।

\medskip\noindent\textbf{भाग I --- दीवार पर लगा नक़्शा।}
\begin{enumerate}
  \item दीवार वाला नक़्शा खींचो: \omterm{def:g8:sexual-reproduction:gametes}{युग्मकों} से आरोपण तक की कड़ी, और उस
        पर \cref{prop:g8:contraception:principle} के तीनों हमले के
        निशान लगे हुए।
  \item पहली आगंतुक पूछती है, ``रोज़ की एक गोली भला रोकेगी क्या?''
        मशीनरी से जवाब दो: गोली किसकी भाषा बोलती है, और किससे?
  \item वही आगे पूछती है: ``और अगर कुछ भूल जाऊँ तो?'' ख़तरे की जगह
        नक़्शे पर बताओ --- \omterm{def:g8:puberty:hormones}{हॉर्मोन} का इशारा टूटने पर क्या फिर से चालू
        हो जाता है?
  \item सलाहकार की \omterm{def:g8:contraception:condom}{कंडोम} वाली सिफ़ारिश हर दूसरे तरीक़े के चुनाव के बाद
        भी क्यों बची रहती है? (वही दोहरी ड्यूटी।)
\end{enumerate}

\medskip\noindent\textbf{भाग II --- हालात।}
\begin{enumerate}[resume]
  \item दूसरा आगंतुक, सत्रह साल का और परेशान, यह मानता रहा कि
        सुरक्षित-दिन वाला कोई ऐप बाक़ी तरीक़े ग़ैर-ज़रूरी कर देता है।
        इस मान्यता को नरमी से सुधारो: कौन-से दो जैविक तथ्य कैलेंडर को
        हरा देते हैं?
  \item तीसरी आगंतुक को ऐसा कुछ चाहिए जिसमें ``रोज़ याद न रखना पड़े''।
        अलमारी के दोनों उम्मीदवार और उनकी कड़ियाँ बताओ।
  \item चौथी आगंतुक \omterm{def:g8:contraception:condom}{कंडोम} फटने के एक दिन के भीतर पहुँचती है। उस विकल्प
        का नाम बताओ, उसकी घड़ी, उसकी मुख्य कार्य-प्रणाली --- और वह
        पक्का तरीक़ा भी, जिस पर सलाहकार आगे के लिए बात करेगी।
  \item चारों संक्रमणों के बारे में एक ही आख़िरी पंक्ति लेकर लौटते
        हैं। वह पंक्ति लिखो।
\end{enumerate}

\medskip\noindent\textbf{भाग III --- समीक्षा।}
\begin{enumerate}[resume]
  \item प्रशिक्षु पूछता है कि केंद्र सिर्फ़ तरीक़े थमाने के बजाय यह
        कड़ी सिखाता ही क्यों है।
        \cref{met:g8:contraception:evaluate} के पहले सवाल से जवाब
        दो।
  \item दोपहर के तरीक़ों को इस हिसाब से छाँटो कि हर एक अपने इस्तेमाल
        करने वाले से क्या माँगता है (तरीक़े का चौथा सवाल)।
  \item प्रशिक्षु सोचता है कि नाबालिगों की मुलाक़ातें बनावट से ही
        गोपनीय क्यों हैं। \cref{ex:g8:contraception:help} के तर्क से
        जवाब दो।
  \item दोपहर का समापन एक वाक्य में करो: डिब्बों से आगे, यह केंद्र
        असल में बाँटता क्या है।
\end{enumerate}
\end{problem}
